\documentclass[10pt]{article}
\usepackage[T1]{fontenc}
\usepackage[utf8]{inputenc}
\usepackage[french]{babel}
\usepackage{amsmath,amssymb,mathrsfs,mathpazo}
\usepackage{geometry}
\geometry{a4paper,margin=2cm}
\usepackage{graphicx,array,tabularx,tabularray,colortbl,multirow,multicol,booktabs,enumitem,pifont,ulem,fancybox,diagbox,cancel}
\usepackage{pstricks,pstricks-add,pst-plot,pst-tree,pst-node,pst-text,pst-func,pst-math,pst-bspline,pst-3dplot,pst-eucl,pst-solides3d}
\let\pscalculate\undefined
\usepackage{pst-bezier}
\usepackage{tikz,pgfplots,tkz-tab}
\usetikzlibrary{arrows,arrows.meta,patterns,calc,positioning,trees,decorations.markings,shapes.geometric,shapes.misc,shapes.arrows,intersections,decorations.pathmorphing,decorations.pathreplacing,backgrounds,shadows,babel}
\pgfplotsset{compat=1.18}
\usepgfplotslibrary{fillbetween}
\usepackage[np]{numprint}
\usepackage{listings,makecell,pas-tableur,twoopt}
\let\barre\undefined
\usepackage[pstricks]{tablvar}
\definecolor{darkgreen}{RGB}{0,100,0}
\providecommand{\red}{\color{red}}
\providecommand{\blue}{\color{blue}}
\providecommand{\green}{\color{green}}
\providecommand{\black}{\color{black}}
\providecommand{\white}{\color{white}}
\providecommand{\gray}{\color{gray}}
\providecommand{\cyan}{\color{cyan}}
\providecommand{\magenta}{\color{magenta}}
\providecommand{\yellow}{\color{yellow}}
\providecommand{\violet}{\color{violet}}
\providecommand{\orange}{\color{orange}}
\providecommand{\euro}{\texteuro}
\providecommand{\EUR}{\texteuro}
\providecommand{\ieme}{\textsuperscript{e}}
\providecommand{\degres}{\ensuremath{^\circ}}
\providecommand{\degree}{\ensuremath{^\circ}}
\providecommand{\vv}[1]{\overrightarrow{#1}}
\providecommand{\Coord}[1]{#1}
\providecommand{\decosix}{}
\providecommand{\starredbullet}{\ensuremath{\bullet}}
\pagestyle{empty}
\setlength{\parindent}{0pt}
\usepackage{adjustbox}\DeclareUnicodeCharacter{2236}{:}\DeclareUnicodeCharacter{2190}{\ensuremath{\leftarrow}}
\providecommand{\R}{\mathbb{R}}
\providecommand{\N}{\mathbb{N}}
\providecommand{\Z}{\mathbb{Z}}
\providecommand{\Q}{\mathbb{Q}}
\providecommand{\D}{\mathbb{D}}
\providecommand{\C}{\mathbb{C}}
\providecommand{\e}{{\mathrm{e}}}
\providecommand{\vect}[1]{\overrightarrow{#1}}
\providecommand{\vectt}[1]{\overrightarrow{\mathrm{#1}}}
\providecommand{\vvt}[1]{\overrightarrow{\mathrm{#1}}}
\providecommand{\np}[2][]{{#2~#1}}
\providecommand{\vv}[1]{{\overrightarrow{#1}}}
\renewcommand{\d}{{\mathrm{d}}}
\providecommand{\degres}{{^\circ}}
\providecommand{\degre}{{^\circ}}
\providecommand{\blue}{\color{blue}}
\providecommand{\red}{\color{red}}
\providecommand{\enskip}{\quad}
\providecommand{\llbracket}{\lbrack\!\lbrack}
\providecommand{\rrbracket}{\rbrack\!\rbrack}
\providecommand{\euro}{{\text{€}}}
\providecommand{\no}{{\text{n°}}}
\providecommand{\interval}[2]{{[#1;#2]}}
\providecommand{\intervFO}[2]{{[#1;#2[}}
\providecommand{\dotfill}{\ldots}

\providecommand{\ds}{\displaystyle}
\providecommand{\textding}[1]{#1}
\providecommand{\Oij}{\(\mathrm{O};\vec{\imath},\vec{\jmath}\)}
\providecommand{\Oijk}{\(\mathrm{O};\vec{\imath},\vec{\jmath},\vec{k}\)}
\providecommand{\Ouv}{\(\mathrm{O};\vec{u},\vec{v}\)}
\providecommand{\intervFF}[2]{\ensuremath{[#1;#2]}}
\renewcommand{\euro}{\ifmmode\text{\texteuro}\else\texteuro\fi}
\DeclareUnicodeCharacter{00B0}{\ifmmode^{\circ}\else\textdegree\fi}
\DeclareUnicodeCharacter{20AC}{\ifmmode\text{\texteuro}\else\texteuro\fi}
% Énoncés : Éducation nationale. Transcription et corrigés : APMEP, Denis Vergès et contributeurs. Découpage et classement : CoopMaths, Sébastien Lozano, Éric Elter et contributeurs. Adaptation de la présentation : Maths Explorer.
% CC BY-SA : https://coopmaths.fr/www/about/
% https://coopmaths.fr/alea/static/e3c/2020/tex/e3c_2020_00_sujet24_2.tex
\begin{document}
\noindent\textit{APMEP / CoopMaths. CC BY-SA. Présentation adaptée par Maths Explorer.}\par\medskip
\medskip

On considère une fonction $f$ définie et dérivable sur l’intervalle $[-4~;~2]$.

La fonction dérivée de $f$ est notée $f'$.

Dans le repère orthonormé ci-dessous, la courbe $\mathcal{C}$ est la courbe représentative de $f$ sur l’intervalle $[-4~;~2]$.

Le point A est le point de la courbe $\mathcal{C}$ d’abscisse $- 1$.

La droite $\mathcal{T}$ est la tangente à la courbe $\mathcal{C}$ en A.

\begin{center}
\psset{xunit=1.8cm,yunit=1.8cm,labelFontSize=\scriptstyle,showorigin=false,comma=true}
\begin{pspicture}(-5,-2.1)(3,3.15)
\multido{\n=-4.3+0.1}{69}{\psline[linewidth=0.3pt,linecolor=lightgray](\n,-2.1)(\n,3)}
\multido{\n=-2.1+0.1}{51}{\psline[linewidth=0.3pt,linecolor=lightgray](-4.3,\n)(2.5,\n)}
\multido{\n=-4+0.5}{14}{\psline[linewidth=0.55pt](\n,-2.1)(\n,3)}
\multido{\n=-2+0.5}{11}{\psline[linewidth=0.55pt](-4.3,\n)(2.5,\n)}
\psaxes[linewidth=0.85pt,,Dy=0.5,]{->}(0,0)(-4.3,-2.1)(2.5,3.1)
 \def\Func{2.71828 x exp x x neg 2.5 add mul 1 sub mul}
 \psplot[plotpoints=2000,linewidth=1.25pt,linecolor=blue]{-4}{2}{\Func}
 \psplot[plotpoints=2000,linewidth=0.75pt,linecolor=green]{-4.3}{2.5}{1.6654575 neg}
\uput[ur](-4,-0.85){\Large\blue $\displaystyle \mathcal{C}$}
\uput[dl](-4,-1.6){\Large\textcolor{green} {$\displaystyle \mathcal{T}$}}
\psdot[dotstyle=+,dotscale=1.8,linecolor=red](-1,-1.6654575)
\uput[d](-1,-1.6654575){A}
\end{pspicture}
\end{center}

\begin{enumerate}
\item Par lecture graphique, donner la valeur de $f'(-1)$.
\item Résoudre, graphiquement, l’inéquation $f'(x)\leqslant 0$.
\end{enumerate}

On admet que la fonction $f$ est définie sur $[-4~;~2]$ par $f(x)=(-x^2 + 2,5x-1)\e^x$.
\begin{enumerate}[resume]
\item Vérifier que, pour tout réel $x$ de l’intervalle $[-4~;~2]$,

\[f'(x)= (- x^2 + 0,5x + 1,5)\e^x.\]

\item Étudier le signe de la fonction $f'$ sur l’intervalle $[-4~;~2]$.
\item En déduire les variations de $f$ sur l’intervalle $[-4~;~2]$.
\end{enumerate}

\vspace{0,5cm}

\end{document}
